\documentclass[11pt,a4paper]{article}
\usepackage[utf8]{inputenc}
\usepackage[T1]{fontenc}
\usepackage{amsmath,amssymb,amsthm}
\usepackage{mathtools}
\usepackage{physics}
\usepackage{hyperref}
\usepackage{enumitem}
\usepackage{cleveref}

\title{Субквантовая теория: ядро}
\author{На основе Теории Мод}
\date{6 июля 2026}

\newtheorem{definition}{Определение}[section]
\newtheorem{axiom}{Аксиома}[section]
\newtheorem{theorem}{Теорема}[section]
\newtheorem{lemma}{Лемма}[section]
\newtheorem{corollary}{Следствие}[section]
\newtheorem{hypothesis}{Гипотеза}[section]
\newtheorem{prediction}{Предсказание}[section]
\newtheorem{remark}{Замечание}[section]

\theoremstyle{definition}
\newtheorem{example}{Пример}[section]

\begin{document}
\maketitle
\tableofcontents
\newpage

% =================================================================
\section{Часть 0: Ядро}
% =================================================================

\subsection{Пространство состояний}

Пусть сеть содержит $N$ частиц, пронумерованных $1, 2, \dots, N$.

\begin{definition}[Конфигурация сети]
Конфигурация $\Sigma(t)$ в момент $t$ задаётся:
\begin{itemize}
    \item Множеством активных частиц $P(t) \subseteq \{1, \dots, N\}$.
    \item Для каждой упорядоченной пары $(A, B)$, $A \neq B$, $A, B \in P(t)$:
    \begin{itemize}
        \item Кратчайшим расстоянием $d_{AB}(t)$.
        \item Поперечным спектром вдоль кратчайшего пути: $F_{A \to B}(k, t)$, где $k = 0, 1, \dots, d_{AB}(t)$.
    \end{itemize}
    \item Для каждой частицы $A \in P(t)$:
    \begin{itemize}
        \item Собственной модой (causal history): $F_{A \to A}(k, t)$, где $k = 0, 1, \dots, d_{AA}(t)$, при этом $d_{AA}(t) \geq 1$.
    \end{itemize}
\end{itemize}

Для частиц, ушедших в пассивный лист ($\text{active} = \text{false}$), спектры $F_{A \to B}$ для всех $B$ продолжают существовать и эволюционировать, но только для $k > d_{AB}$ (хвосты). Наблюдаемые участки ($k \leq d_{AB}$) для пассивных частиц заморожены на момент деактивации.
\end{definition}

\begin{definition}[Гильбертово пространство одной моды]
Для каждой пары $(A, B)$ определён комплексный вектор:
\begin{equation}
    \psi_{AB}(t) = (F_{A \to B}(0, t), F_{A \to B}(1, t), \dots, F_{A \to B}(d_{AB}(t), t))
\end{equation}
принадлежащий пространству $\mathcal{H}_{AB}(t) = \mathbb{C}^{d_{AB}(t)+1}$.
\end{definition}

\begin{definition}[Полное состояние]
Состояние сети:
\begin{equation}
    \Psi(t) = \{\psi_{AB}(t) \mid A \in P(t), B \in P(t), \text{ путь } A \to B \text{ существует}\}
\end{equation}
лежит в прямом произведении соответствующих пространств.
\end{definition}

% -----------------------------------------------------------------
\subsection{Локальные величины}

\begin{definition}[Корреляция]
\begin{equation}
    C_{AB}(t) = \sum_{k=0}^{d_{AB}} F_{A \to B}(k, t) \cdot F_{B \to A}(d_{AB} - k, t)
\end{equation}
\end{definition}

\begin{definition}[Энергия пары]
\begin{equation}
    E_{AB}(t) = |C_{AB}(t)|
\end{equation}
\end{definition}

\begin{definition}[Полная энергия сети]
\begin{equation}
    E_{\text{total}}(t) = \sum_{A < B} |C_{AB}(t)|
\end{equation}
\end{definition}

\begin{definition}[Гравитационная масса]
\begin{equation}
    m_{\text{grav}}^A(t) = \Lambda \cdot \operatorname{sign}\!\left(\operatorname{Re} \sum_{k=0}^{d_{AA}} F_{A \to A}(k)\right) \cdot \left[ \sum_{k=0}^{d_{AA}} |F_{A \to A}(k, t)|^2 \right]^{1/2}
\end{equation}
где $\Lambda$ --- фундаментальная константа размерности массы. Знак суммы отражает общую фазу causal history: положительная сумма $\to$ обычная материя, отрицательная $\to$ инвертированная.
\end{definition}

\begin{definition}[Ортогональность]
\begin{equation}
    O_A(t) = 1 - \frac{m_{\text{grav}}^A(t)}{m_p}
\end{equation}
где $m_p$ --- масса протона.
\end{definition}

\begin{definition}[Инертная масса]
\begin{equation}
    m_{\text{inert}}^A(t) = \frac{S_A(t)}{O_A(t)}
\end{equation}
где $S_A(t) = \sum_{k=0}^{K_0-1} |F_{A \to A}(k, t)|^2$ --- площадь локального участка собственной моды,
$K_0 = \min\{k > 0 \mid \exists B \neq A: \text{в causal history узла } k \text{ записана свёртка с } B\}$.
\end{definition}

\begin{definition}[Энергия частицы]
\begin{equation}
    E_A(t) = \frac{1}{|H_A|} \sum_{B \in H_A} |C_{AB}(t)|
\end{equation}
где $H_A$ --- причинный горизонт $A$.
\end{definition}

\begin{corollary}[Сохранение $E + m + O$]
Для частиц одного семейства:
\begin{equation}
    E_A(t) + m_{\text{inert}}^A(t) + O_A(t) = \text{const}
\end{equation}
\end{corollary}

% -----------------------------------------------------------------
\subsection{Калибровочные условия}

\begin{axiom}[Противофазность причинных пар]
Для любых $A, B$, связанных причинно (т.е. существует путь $A \to B$, по которому шёл градиентный спуск):
\begin{equation}
    F_{A \to B}(0, t) = -F_{B \to A}(0, t)
\end{equation}
Данная аксиома выполняется строго в равновесии (при $t \to \infty$ вдоль причинной цепи). В моменты рождения и исчезновения калибровка может временно нарушаться и восстанавливается градиентным спуском за время порядка $1/\alpha$.
\end{axiom}

\begin{corollary}[Индефинитная метрика]
\begin{equation}
    \eta_{AB} = \operatorname{sign}(F_{A \to B}(0) \cdot F_{B \to A}(0)) = -1 \quad \text{для всех причинных пар в равновесии}
\end{equation}
\end{corollary}

\begin{axiom}[Глобальная компенсация произведения]
Для любых $A \neq B$:
\begin{equation}
    \sum_{k=0}^{\infty} F_{A \to B}(k, t) \cdot F_{B \to A}(k, t) = 0
\end{equation}
\end{axiom}

% -----------------------------------------------------------------
\subsection{Гравитационное дальнодействие}

\begin{definition}[Гравитационный потенциал]
\begin{equation}
    V_{AB}(t) = G \cdot \sum_{k = d_{AB}+1}^{\infty} \frac{F_{A \to B}(k, t) \cdot F_{B \to A}(k, t)}{k^2}
\end{equation}
где $G$ --- гравитационная постоянная. Хвосты для $k > d_{AB}$ --- проекции causal history остальной Вселенной.

Индекс $k$ вдоль моды связан с физическим расстоянием $r$ через масштабный фактор: $r = \delta \cdot k$, где $\delta$ --- фундаментальная длина (порядка планковской). Формула $V_{AB}$ в непрерывном пределе переходит в $V \propto 1/r$ для трёхмерного пространства-времени.
\end{definition}

\begin{corollary}[Связь локального и дальнего]
Из Аксиомы глобальной компенсации:
\begin{equation}
    \sum_{k=0}^{d_{AB}} F_{A \to B}(k) \cdot F_{B \to A}(k) + \frac{1}{G} \sum_{k > d_{AB}} k^2 \cdot V_{AB}(k) = 0
\end{equation}
Локальное усиление корреляции уменьшает гравитационный хвост и наоборот.
\end{corollary}

% -----------------------------------------------------------------
\subsection{Динамика: оператор шага}

\begin{definition}[Шаг эволюции]
$\Sigma(t+1) = \mathcal{U}[\Sigma(t)]$, где $\mathcal{U} = \mathcal{U}_{\text{spread}} \circ \mathcal{U}_{\text{birth}} \circ \mathcal{U}_{\text{death}} \circ \mathcal{U}_{\text{history}}$.

Порядок применения:
\begin{enumerate}
    \item $\mathcal{U}_{\text{history}}$ (градиентный спуск)
    \item $\mathcal{U}_{\text{birth}}$ (проверка рождений)
    \item $\mathcal{U}_{\text{death}}$ (проверка исчезновений)
    \item $\mathcal{U}_{\text{spread}}$ (рост горизонта)
\end{enumerate}
\end{definition}

\textbf{$\mathcal{U}_{\text{spread}}$ (Рост горизонта).}
Для всех $A \in P(t)$:
\begin{equation}
    d_{AA} \to d_{AA} + 1, \quad H_A \text{ расширяется на } 1
\end{equation}
Хвосты, бывшие при $k > d_{AB}$, теперь при $k = d_{AB}+1$ попадают под действие $\mathcal{U}_{\text{history}}$ на следующем шаге.

При каждом шаге $\mathcal{U}_{\text{spread}}$ оператор сдвига $T$ применяется ко всем модам: $F(k) \to F(k+1)$ для $k \geq 0$, а на освободившееся место $k=0$ записывается новое значение из взаимодействий текущего шага.

\textbf{$\mathcal{U}_{\text{birth}}$ (Рождение).}
Для каждой пары $(A, B)$ в контакте ($d_{AB} = 0$):
\begin{itemize}
    \item Если $|C_{AB}(t)| > \theta$:
    \begin{itemize}
        \item С вероятностью $p$:
        \begin{itemize}
            \item Создать новую частицу $C$.
            \item $P(t+1) = P(t) \cup \{C\}$.
            \item $d_{AC} = 0$, $d_{CB} = 0$ ($C$ помещается между $A$ и $B$).
            \item Инициализировать спектры из уравнения баланса:
            \begin{equation}
                F_{A \to C}(0) \cdot F_{C \to A}(0) + F_{C \to B}(0) \cdot F_{B \to C}(0) = |C_{AB}| - |C_{AB}^{\text{new}}|
            \end{equation}
            с выбором значений, минимизирующих локальную нескомпенсированность.

            Уравнение баланса --- точное равенство, следующее из сохранения causal history при рождении. Полная causal history сети не меняется при рождении частицы; она перераспределяется между старыми и новыми модами. Сумма $|C_{AC}| + |C_{CB}| + |C_{AB}^{\text{new}}| = |C_{AB}^{\text{old}}|$ --- точное равенство.
            \item $d_{CC} = 1$, $F_{C \to C}(0) = F_{C \to C}(1) = 0$.
            \item Удлинить causal history $A$ и $B$ на $1$:
            \begin{equation}
                d_{AA} \mathrel{+}= 1, \quad d_{BB} \mathrel{+}= 1
            \end{equation}
            В новый узел записать свёртку $F_{A \to B}(0) \cdot F_{B \to A}(0)$ в момент рождения.
        \end{itemize}
    \end{itemize}
\end{itemize}

\textbf{$\mathcal{U}_{\text{death}}$ (Исчезновение).}
Для каждой тройки $A$--$C$--$B$, где $d_{AC} = d_{CB} = 0$, $d_{AB} = 1$:
\begin{itemize}
    \item Если $|C_{AB}(t)| < \theta$:
    \begin{itemize}
        \item Удалить $C$ из $P$ (causal history $C$ уходит в пассивный лист).
        \item Установить $d_{AB} = 0$ ($A$ и $B$ сближаются).
        \item Новый спектр $F_{A \to B}^{\text{new}}(k)$ --- свёртка старых $F_{A \to C}$ и $F_{C \to B}$:
        \begin{equation}
            F_{A \to B}^{\text{new}}(k) = e^{i\phi} \sum_{i=0}^{k} F_{A \to C}(i) \cdot F_{C \to B}(k-i), \quad k = 0, \dots, d_{AB}^{\text{new}}
        \end{equation}
        Фаза $\phi$ определяется из условия минимизации энергии новой конфигурации. При симметричном исчезновении $\phi = 0$. При несимметричном $\phi$ выбирается так, чтобы минимизировать $|C_{AB}^{\text{new}}|$ после свёртки.
        \item Пассивный лист $C$ продолжает существовать и эволюционировать по правилам Определения~0.5.2.
    \end{itemize}
\end{itemize}

\textbf{$\mathcal{U}_{\text{history}}$ (Градиентный спуск).}
Для всех $A, B \in P(t)$:
\begin{itemize}
    \item Вычислить $\partial E_{\text{total}} / \partial F_{A \to B}^*(k)$ для всех $k \leq d_{AB}$.
    \item Обновить:
    \begin{equation}
        F_{A \to B}(k, t+1) = F_{A \to B}(k, t) - \alpha \cdot \frac{\partial E_{\text{total}}}{\partial F_{A \to B}^*(k)}
    \end{equation}
    где $\alpha > 0$ --- скорость обучения (фундаментальная постоянная).
\end{itemize}

Хвосты ($k > d_{AB}$) не обновляются градиентным спуском. Они меняются только при росте горизонта, когда $k$, бывшее $> d$, становится $\leq d$ и попадает под действие $\mathcal{U}_{\text{history}}$.

\begin{definition}[Эволюция пассивного листа]
Для частицы $C$, ушедшей в пассивный лист в момент $t_{\text{death}}$:
\begin{itemize}
    \item Её causal history продолжает расти: $d_{CC}$ растёт на $1$ каждый шаг.
    \item В новый узел записывается свёртка хвостов: $F_{C \to C}(d_{\text{new}}) = \sum_{B \in H_C} F_{C \to B}^{\text{tail}}(d_{\text{new}}) \cdot F_{B \to C}^{\text{tail}}(d_{\text{new}})$.
    \item Внешние моды $C$ для активных частиц $B$ существуют только как хвосты ($k > d_{CB}(t_{\text{death}})$).
    \item Гравитационная масса пассивного листа растёт монотонно.
    \item Ортогональность пассивного листа заморожена в значении на момент деактивации.
    \item Пассивный лист может быть реактивирован, если корреляция с активной частицей превысит $\theta$.
\end{itemize}
\end{definition}

% -----------------------------------------------------------------
\subsection{Простейший пример: 2 частицы}

\begin{example}[Начальное состояние $t=0$]
Две частицы $A, B$. $d_{AB} = 0$ (в контакте).
Собственные моды: $d_{AA} = 1$, $d_{BB} = 1$, $F_{A \to A}(0) = 0$, $F_{B \to B}(0) = 0$.

Спектры контакта:
\begin{equation}
    F_{A \to B}(0) = a, \quad F_{B \to A}(0) = -a \quad \text{(калибровка)}
\end{equation}
где $a > 0$.

Вычисляем:
\begin{align}
    C_{AB} &= a \cdot (-a) = -a^2 \\
    |C_{AB}| &= a^2 \\
    E_{\text{total}} &= a^2 \\
    m_{\text{grav}}^A &= \Lambda \cdot \operatorname{sign}(0) \cdot |0| = 0 \\
    O_A &= 1 \\
    m_{\text{inert}}^A &= 0
\end{align}

Интерпретация: два безмассовых кванта в контакте, встречные проекции противофазны, корреляция отрицательна (притяжение).
\end{example}

\begin{example}[Шаг 1: Градиентный спуск]
\begin{align}
    \frac{\partial E}{\partial F_{A \to B}(0)} &= \operatorname{sign}(C_{AB}) \cdot F_{B \to A}(0) = (-1) \cdot (-a) = a \\
    \frac{\partial E}{\partial F_{B \to A}(0)} &= \operatorname{sign}(C_{AB}) \cdot F_{A \to B}(0) = (-1) \cdot a = -a
\end{align}

Обновление:
\begin{align}
    F_{A \to B}(0, 1) &= a - \alpha a = a(1 - \alpha) \\
    F_{B \to A}(0, 1) &= -a - \alpha(-a) = -a(1 - \alpha)
\end{align}

При $\alpha = 1$: $F \to 0$. Система релаксирует к нулевой корреляции. Случай $\alpha = 1$ приведён для иллюстрации и соответствует бесконечному темпу спуска; физически реализуется режим $\alpha \ll 1$.

При $\alpha < 1$: амплитуда убывает, корреляция $|C| = a^2(1-\alpha)^2$ уменьшается.
Рождения не происходит, если $a^2(1-\alpha)^2 \leq \theta$.
\end{example}

\newpage
% =================================================================
\section{Часть 1: Три частицы --- рождение и исчезновение посредника}
% =================================================================

\subsection{Начальная конфигурация}

Три частицы: $A, B, C$.

\textbf{$t = 0$:}
\begin{itemize}
    \item $d_{AB} = 1$ ($A$ и $B$ разделены посредником $C$)
    \item $d_{AC} = 0$, $d_{CB} = 0$ ($C$ в контакте с обоими)
    \item $d_{AA} = 1$, $d_{BB} = 1$, $d_{CC} = 2$
\end{itemize}

Спектры:
\begin{align}
    F_{A \to C}(0) &= a, \quad F_{C \to A}(0) = -a \\
    F_{C \to B}(0) &= b, \quad F_{B \to C}(0) = -b
\end{align}

Собственная мода $C$: $F_{C \to C}(0) = \gamma_0$, $F_{C \to C}(1) = \gamma_1$.

\begin{definition}[Спектр вдоль непрямого пути]
Для пути $\Gamma = A \to X_1 \to X_2 \to \dots \to X_d \to B$:
\begin{equation}
    F_{A \to B}(0) = F_{A \to X_1}(0)
\end{equation}
Для $k \geq 1$:
\begin{equation}
    F_{A \to B}(k) = \frac{\prod_{i=1}^{k} F_{X_{i-1} \to X_i}(0)}{N_k}, \quad N_k = \sum_{i=1}^{k} |F_{X_{i-1} \to X_i}(0)|
\end{equation}
где $N_k$ --- нормировка, предотвращающая экспоненциальный рост.

В случае $d=1$ (один посредник $C$):
\begin{equation}
    F_{A \to B}(0) = a, \quad F_{A \to B}(1) = b
\end{equation}
\end{definition}

Калибровка требует $F_{A \to B}(0) = -F_{B \to A}(0) \Rightarrow a = b$. Это упрощение для симметричного случая. В общем случае $a \neq b$, и калибровка даёт $a = -F_{B \to A}(0)$, $b = -F_{A \to B}(0)$, что всегда выполнимо.

Принимаем $a = b$. Тогда:
\begin{equation}
    F_{A \to B} = (a, a), \quad F_{B \to A} = (-a, -a)
\end{equation}

\subsection{Корреляции и энергия}

\begin{align}
    C_{AB} &= \sum_{k=0}^{1} F_{A \to B}(k) \cdot F_{B \to A}(1-k) \notag\\
           &= a \cdot (-a) + a \cdot (-a) = -2a^2 \\
    |C_{AB}| &= 2a^2 \\
    C_{AC} &= -a^2, \quad C_{CB} = -a^2 \\
    E_{\text{total}} &= 2a^2 + a^2 + a^2 = 4a^2
\end{align}

\subsection{Сценарий 1: Исчезновение $C$}

Условие: $|C_{AB}| < \theta$.

После градиентного спуска корреляция $C_{AB}$ падает. При срабатывании порога $C$ исчезает:

\begin{enumerate}
    \item $C$ удаляется из $P$ ($\text{active} = \text{false}$, causal history $C$ уходит в пассивный лист).
    \item $A$ и $B$ входят в прямой контакт: $d_{AB} = 0$.
    \item Новый спектр --- свёртка:
    \begin{equation}
        F_{A \to B}^{\text{new}}(k) = e^{i\phi} \sum_{i=0}^{k} F_{A \to C}^{\text{old}}(i) \cdot F_{C \to B}^{\text{old}}(k-i)
    \end{equation}
\end{enumerate}

Фаза $\phi$ определяется из условия минимизации энергии новой конфигурации. При симметричном исчезновении $\phi = 0$. При несимметричном $\phi$ выбирается так, чтобы минимизировать $|C_{AB}^{\text{new}}|$ после свёртки.

\begin{remark}[Динамическая калибровка]
Калибровка $F_{A \to B}(0) = -F_{B \to A}(0)$ выполняется в равновесии ($t \to \infty$ вдоль причинной цепи), а в моменты рождения/исчезновения может временно нарушаться. Градиентный спуск восстанавливает калибровку за время порядка $1/\alpha$.
\end{remark}

\subsection{Сценарий 2: Рождение $D$ между $A$ и $B$}

Начало: $C$ исчезла, $A$ и $B$ в контакте, $d_{AB} = 0$.

$F_{A \to B}(0) = u$, $F_{B \to A}(0) = -u$ (калибровка восстановлена).

Условие рождения: $|C_{AB}| = u^2 > \theta$.

Рождается $D$ между $A$ и $B$. Симметричное решение:
\begin{equation}
    F_{A \to D}(0) = v, \quad F_{D \to A}(0) = -v, \quad F_{D \to B}(0) = v, \quad F_{B \to D}(0) = -v
\end{equation}

Уравнение баланса рождения:
\begin{equation}
    |C_{AB}^{\text{new}}| = |C_{AB}| - (|C_{AD}| + |C_{DB}|)
\end{equation}
Сумма $|C_{AD}| + |C_{DB}| + |C_{AB}^{\text{new}}| = |C_{AB}^{\text{old}}|$ --- точное равенство.

$|C_{AD}| = v^2$, $|C_{DB}| = v^2$.
$|C_{AB}^{\text{new}}| = u^2 - 2v^2$.

Выбираем $v = u/2$: $|C_{AB}^{\text{new}}| = u^2/2$. Рождение уменьшило корреляцию $A$--$B$ вдвое, перераспределив её на новые пары.

\subsection{Градиентный спуск в системе трёх частиц}

Для конфигурации $A$--$C$--$B$:
\begin{equation}
    E_{\text{total}} = |F_{A \to C} \cdot F_{C \to A}| + |F_{C \to B} \cdot F_{B \to C}| + |F_{A \to C} \cdot F_{C \to B} + F_{C \to B} \cdot F_{C \to A}|
\end{equation}

В симметричной точке:
\begin{equation}
    \frac{\partial E}{\partial a} = 8a, \quad a(t+1) = a(t)(1 - 8\alpha)
\end{equation}

Условие устойчивости: $\alpha < 1/8$ (для симметричного случая). В общем случае $\alpha_{\max} = 1 / (2 + 2\sqrt{2}) \approx 0.207$.

\newpage
% =================================================================
\section{Часть 2: Спектральное представление}
% =================================================================

\subsection{Мода как функция на дискретной оси}

Для пары $(A, B)$ определён спектр $F_{A \to B} : \mathbb{N}_0 \to \mathbb{C}$, образующий пространство $\ell^2(\mathbb{N}_0)$ с нормой:
\begin{equation}
    \|F_{A \to B}\|^2 = \sum_{k=0}^{\infty} |F_{A \to B}(k)|^2
\end{equation}

\subsection{Оператор сдвига $T$}

\begin{definition}[Оператор сдвига вдоль моды]
\begin{equation}
    (T F)(k) = F(k+1), \quad k \geq 0
\end{equation}
\end{definition}

Свойства:
\begin{itemize}
    \item $T$ --- изометрия на своём образе: $\|T F\| \leq \|F\|$, с равенством если $F(0) = 0$.
    \item $T^\dagger$ действует как $(T^\dagger F)(0) = 0$, $(T^\dagger F)(k) = F(k-1)$ для $k \geq 1$.
    \item $T^\dagger T \neq I$ (не унитарен --- есть потеря на границе $k=0$).
\end{itemize}

Физический смысл: потеря на границе $k=0$ соответствует тому, что causal history необратима: информация, ушедшая в $k=0$, не может быть восстановлена. При каждом шаге $\mathcal{U}_{\text{spread}}$ оператор $T$ применяется ко всем модам: $F(k) \to F(k+1)$, а на $k=0$ записывается новое значение.

\subsection{Собственные моды оператора эволюции}

Рассмотрим собственную моду частицы $A$: $\psi_A(k) = F_{A \to A}(k)$.

В непрерывном пределе градиентный спуск даёт:
\begin{equation}
    \pdv{\psi_A}{t} = -\alpha \frac{\delta E}{\delta \psi_A^*}
\end{equation}

Вблизи равновесия линеаризуем:
\begin{equation}
    \pdv{\psi_A}{t} = -H_A \psi_A
\end{equation}
где $H_A$ --- эрмитов оператор.

\begin{definition}[Спектральное разложение]
Собственные функции $H_A$:
\begin{equation}
    H_A \phi_n = \lambda_n \phi_n, \quad \lambda_n \geq 0
\end{equation}
Общее решение:
\begin{equation}
    \psi_A(t) = \sum_n c_n e^{-\lambda_n t} \phi_n
\end{equation}
\end{definition}

\subsection{Дисперсионное соотношение}

\begin{hypothesis}[Дисперсионное соотношение]
Для моды между $A$ и $B$:
\begin{equation}
    \lambda_n \propto \frac{n}{d_{AB}}
\end{equation}
\end{hypothesis}

Данная гипотеза --- линейное приближение. Для реальных мод с плато и пиками распределение может иметь сгущения и разрежения, что не меняет качественных выводов.

\subsection{Представление корреляции через собственные моды}

\begin{align}
    \psi_{A \to B} &= \sum_n a_n \phi_n^{AB}, \quad \psi_{B \to A} = \sum_m b_m \phi_m^{BA} \\
    C_{AB} &= \sum_{n,m} a_n b_m \cdot \left( \sum_{k=0}^{d} \phi_n^{AB}(k) \cdot \phi_m^{BA}(d-k) \right)
\end{align}

Для ортонормированных $\phi$ с экспоненциальным убыванием $\phi_n(k) \propto e^{-\kappa_n k}$:
\begin{equation}
    C_{AB} \propto e^{-\kappa_{\min} \cdot d_{AB}}
\end{equation}

\begin{corollary}[Конфайнмент]
Корреляция экспоненциально подавлена с расстоянием. Константа $\kappa_{\min}$ зависит от типа взаимодействия: для сильного --- велика (конфайнмент на масштабе ферми), для электромагнитного --- мала или ноль (дальнодействие).
\end{corollary}

\subsection{Гравитационные моды}

Гравитационное дальнодействие формируется модами с $\lambda_n \approx 0$. Для них $\phi_n(k) \approx \text{const}$ при $k \to \infty$. В трёхмерном пространстве-времени это даёт $V \propto 1/r$.

\subsection{Квантование спектра}

\begin{definition}[Квантовые числа как топологические инварианты]
Каждая собственная мода $\phi_n$ имеет узловую структуру (число нулей на $[0, d_{AB}]$). Число нулей --- топологический инвариант.
\end{definition}

\begin{corollary}
Квантовые числа (заряд, цвет, аромат) --- числа нулей соответствующих собственных мод частицы. Спин --- узловая структура собственной моды $\phi_0$.
\end{corollary}

\subsection{Фермионы и бозоны}

\begin{itemize}
    \item \textbf{Бозон:} Собственная мода основного тона имеет чётное число нулей. Один оборот ($360^\circ$) возвращает фазу к исходной.
    \item \textbf{Фермион:} Собственная мода основного тона имеет нечётное число нулей. Один оборот меняет знак, два оборота ($720^\circ$) возвращают фазу.
\end{itemize}

Предсказание: для электрона (спин $1/2$) основная мода имеет ровно один ноль. Для фотона (спин $1$) --- ноль или два. Может быть проверено через анализ шумовых паттернов.

\subsection{Информационный поток вдоль моды}

\begin{definition}[Ток информации]
\begin{equation}
    J_{A \to B}(k) = \operatorname{Im}\left[ F_{A \to B}^*(k) \cdot F_{A \to B}(k+1) \right]
\end{equation}
\end{definition}

Сохранение:
\begin{equation}
    \sum_{k=0}^{d-1} J_{A \to B}(k) = \operatorname{Im}\left[ F_{A \to B}^*(0) \cdot F_{A \to B}(d) \right]
\end{equation}

Если $F(0)$ и $F(d)$ синфазны --- ток сохраняется. Если нет --- информация утекает в хвост (гравитационное излучение).

\newpage
% =================================================================
\section{Часть 3: Вывод правила Борна}
% =================================================================

\subsection{Наблюдатель и его горизонт}

\begin{definition}[Горизонт наблюдателя]
Наблюдатель $O$ --- выделенная частица сети. Его причинный горизонт $H_O(t)$ конечен. Для любой другой частицы $S$ наблюдатель видит только:
\begin{equation}
    F_{S \to O}^{\text{obs}}(k), \quad k = 0, 1, \dots, d_{SO}
\end{equation}
Хвост $k > d_{SO}$ ненаблюдаем.
\end{definition}

\begin{definition}[Макросостояние]
Макросостояние системы $S$ относительно наблюдателя $O$:
\begin{equation}
    \text{Data}_O(S) = \{ F_{S \to O}(k) \mid k \leq d_{SO} \} \cup \{ F_{S \to B}(k) \mid B \in H_O, k \leq d_{SB} \} \cup \{ F_{S \to S}(k) \mid k \leq \min(d_{SS}, H_O) \}
\end{equation}
Вся остальная информация --- скрытые переменные (хвосты мод).
\end{definition}

\subsection{Ансамбль микросостояний}

\begin{definition}[Ансамбль]
\begin{equation}
    \Omega[\text{Data}_O(S)] = \{ \text{все допустимые продолжения спектров на } k > d_{SO}, \text{ совместимые с наблюдаемыми данными и аксиомами сети} \}
\end{equation}
\end{definition}

\begin{axiom}[Принцип равновероятности]
В отсутствие дополнительной информации все микросостояния в $\Omega$ равновероятны (следствие максимальной энтропии при фиксированных наблюдаемых).
\end{axiom}

\subsection{Амплитуда перехода}

\begin{definition}[Амплитуда перехода]
\begin{equation}
    A(\text{in} \to \text{out}) = \sum_{k=0}^{d_{\text{obs}}} F_{S \to O}^{\text{in}}(k) \cdot F_{O \to S}^{\text{out}}(d_{\text{obs}} - k)
\end{equation}
где $d_{\text{obs}} = \min(d_{\text{in}}, d_{\text{out}}, |H_O|)$.
\end{definition}

\subsection{Усреднение по хвостам}

Разложим полную корреляцию: $C_{\text{full}} = C_{\text{obs}} + C_{\text{tail}}$, где $C_{\text{obs}} = A(\text{in} \to \text{out})$.

\begin{equation}
    |C_{\text{full}}|^2 = |C_{\text{obs}}|^2 + C_{\text{obs}}^* C_{\text{tail}} + C_{\text{obs}} C_{\text{tail}}^* + |C_{\text{tail}}|^2
\end{equation}

\begin{lemma}[Некоррелированность хвостов]
Для макроскопически различных in и out, хвосты $F_{\text{in}}(k)$ и $F_{\text{out}}(k)$ для $k > d_{\text{obs}}$ имеют случайные относительные фазы при усреднении по $\Omega$.

Обоснование: Фазы хвостов случайны в силу Аксиомы глобальной компенсации произведения. Сумма $\sum F \cdot F' = 0$ накладывает одно интегральное ограничение на фазы, но не фиксирует их индивидуально. В термодинамическом пределе $N \to \infty$ фазы хвостов распределены равномерно на $[0, 2\pi)$.
\end{lemma}

\begin{corollary}
\begin{equation}
    \langle C_{\text{tail}} \rangle_\Omega = 0, \quad \langle C_{\text{obs}}^* C_{\text{tail}} \rangle_\Omega = 0, \quad \langle |C_{\text{tail}}|^2 \rangle_\Omega = \Sigma^2
\end{equation}
\end{corollary}

\subsection{Вывод правила Борна}

\begin{align}
    P(\text{in} \to \text{out}) &= \frac{\langle |C_{\text{full}}|^2 \rangle_\Omega}{\mathcal{N}} = \frac{|A|^2 + \Sigma^2}{\mathcal{N}} \\
    \mathcal{N} &= \sum_{\text{out}'} |A(\text{in} \to \text{out}')|^2 + M\Sigma^2
\end{align}

В пределе большого горизонта:
\begin{equation}
    \boxed{P(\text{in} \to \text{out}) = \frac{|A(\text{in} \to \text{out})|^2}{\sum_{\text{out}'} |A(\text{in} \to \text{out}')|^2}}
\end{equation}

Это правило Борна.

\subsection{Инвариантность относительно глобальной фазы}

При умножении всех $F_{S \to O}(k)$ на $e^{i\phi}$:
\begin{equation}
    A \to e^{i(\phi_{\text{in}} - \phi_{\text{out}})} A, \quad |A|^2 \text{ не меняется}
\end{equation}

\subsection{Интерференция}

Если система может перейти двумя путями: $A_{\text{total}} = A_1 + A_2$.
\begin{equation}
    P = |A_1 + A_2|^2 = |A_1|^2 + |A_2|^2 + 2 \operatorname{Re}(A_1^* A_2)
\end{equation}
Интерференционный член возникает из-за общих хвостов.

\subsection{Измерение как синхронизация}

\begin{definition}[Измерение]
Измерение величины $Q$ у системы $S$ прибором $A$ --- установление прямого контакта ($d_{AS} = 0$) и градиентный спуск до порога, при котором causal history $A$ и $S$ синхронизируются.
\end{definition}

После измерения:
\begin{itemize}
    \item Наблюдаемый участок спектра $S$ ($k \leq d_{AS}$) синхронизируется с прибором $A$.
    \item Хвосты $S$ ($k > d_{AS}$) не меняются мгновенно. Они продолжают эволюционировать по обычным правилам.
    \item Унитарность полной системы сохраняется: коллапс --- только синхронизация наблюдаемых участков.
    \item С точки зрения наблюдателя это коллапс волновой функции. С точки зрения полной сети --- обратимая запись данных в causal history.
\end{itemize}

\newpage
% =================================================================
\section{Часть 4: Вывод уравнения Шрёдингера}
% =================================================================

\subsection{Непрерывный предел}

Введём непрерывные переменные:
\begin{equation}
    t' = \varepsilon \cdot t, \quad x = \delta \cdot k, \quad \psi(x, t') = \lim F(k, t)
\end{equation}
где $\varepsilon, \delta \to 0$, $\delta/\varepsilon = c$ (скорость света).

\subsection{Градиентный спуск в непрерывном пределе}

Дискретный спуск: $F(k, t+1) = F(k, t) - \alpha \cdot \partial E / \partial F^*(k, t)$.

В непрерывном пределе:
\begin{equation}
    \pdv{\psi}{t'} = -\alpha' \frac{\delta E}{\delta \psi^*}, \quad \alpha' = \frac{\alpha}{\varepsilon}
\end{equation}

Энергия сети:
\begin{equation}
    E[\psi] = \int dx \left[ \left| \pdv{\psi}{x} \right|^2 + V(x)|\psi|^2 + \lambda|\psi|^4 + \dots \right]
\end{equation}

\subsection{От диффузии к осцилляциям}

Вблизи равновесия, удерживая квадратичные члены:
\begin{equation}
    \pdv{\psi}{t'} = \alpha' \pdv[2]{\psi}{x} - \alpha' V(x) \psi
\end{equation}

Это уравнение диффузии. Теперь ключевой шаг.

Рассмотрим две моды $A \to B$ и $B \to A$. Калибровка $F_{A \to B}(0) = -F_{B \to A}(0)$ связывает их. Введём переменные:
\begin{equation}
    u = \operatorname{Re}(F_{A \to B}(0)), \quad v = \operatorname{Im}(F_{A \to B}(0))
\end{equation}
Тогда $F_{B \to A}(0) = -u - iv$.

Градиентный спуск по $E = |F_{A \to B}(0) \cdot F_{B \to A}(0)| = u^2 + v^2$ даёт:
\begin{equation}
    \dv{u}{t} \propto -u, \quad \dv{v}{t} \propto -v
\end{equation}

Но калибровка действует на каждом шаге, пересчитывая $F_{B \to A}(0)$ через $F_{A \to B}(0)$. Это вводит обратную связь. В непрерывном пределе:
\begin{equation}
    \dv{u}{t} = -\omega v, \quad \dv{v}{t} = \omega u
\end{equation}
где $\omega \propto \alpha/\varepsilon$. Это осцилляции.

Объединяя $u$ и $v$ в комплексное $\psi = u + iv$:
\begin{equation}
    i \dv{\psi}{t} = \omega \psi
\end{equation}

С учётом потенциала $V(x)$:
\begin{equation}
    \boxed{i \hbar_{\text{eff}} \pdv{\psi}{t} = H_{\text{eff}} \psi}
\end{equation}

Мнимая единица возникает из-за обмена информацией между действительной и мнимой частями спектра, вызванного калибровочными связями. Уравнение Шрёдингера --- эффективное описание градиентного спуска с калибровочными ограничениями.

\subsection{Стационарное уравнение}

Для $\psi(x, t) = \phi(x) e^{-iEt/\hbar}$:
\begin{equation}
    H_{\text{eff}} \phi = E \phi
\end{equation}
$E_n = \hbar \lambda_n$ --- энергия $n$-го состояния.

\subsection{Связь с квантовой механикой}

\begin{center}
\begin{tabular}{c|c}
\textbf{Квантовая механика} & \textbf{Субквантовая теория} \\
\hline
Волновая функция $\psi(x)$ & Наблюдаемый участок спектра $F(k)$, $k \leq d_{\text{obs}}$ \\
Уравнение Шрёдингера & Градиентный спуск с калибровочными связями \\
Правило Борна & Усреднение по ненаблюдаемым хвостам \\
Операторы наблюдаемых & Проекции корреляций на горизонт \\
Коллапс & Запись в causal history \\
Постоянная Планка $\hbar$ & $\varepsilon \cdot \alpha'$ (шаг дискретизации $\times$ темп спуска)
\end{tabular}
\end{center}

$\hbar$ --- эмерджентная константа, выводимая из статистики сети.

\newpage
% =================================================================
\section{Часть 5: Численная модель}
% =================================================================

\subsection{Архитектура}

Модель реализует:
\begin{itemize}
    \item Сеть из $N$ частиц с переменной топологией.
    \item Спектры $F_{A \to B}(k)$ на модах.
    \item Рождение/исчезновение частиц по порогам корреляции.
    \item Градиентный спуск энергии.
    \item Сбор статистики и визуализацию.
\end{itemize}

\subsection{Структуры данных}

\begin{verbatim}
class Particle:
    id: int
    active: bool
    causal_history: list[complex]
    horizon: set[int]

class Mode:
    source: int
    target: int
    distance: int
    spectrum: list[complex]
    tail: list[complex]

class Network:
    particles: dict[int, Particle]
    modes: dict[(int, int), Mode]
    time: int
\end{verbatim}

\subsection{Алгоритм шага эволюции}

\begin{enumerate}
    \item \textbf{Градиентный спуск} ($\mathcal{U}_{\text{history}}$): обновить все $F(k)$ по $F \leftarrow F - \alpha \cdot \partial E/\partial F^*$.
    \item \textbf{Рождение} ($\mathcal{U}_{\text{birth}}$): если $d_{AB} = 0$ и $|C_{AB}| > \theta$, с вероятностью $p$ создать $C$, инициализировать спектры из уравнения баланса.
    \item \textbf{Исчезновение} ($\mathcal{U}_{\text{death}}$): если $A$--$C$--$B$ и $|C_{AB}| < \theta$, удалить $C$ (causal history в пассивный лист), создать прямую моду $A$--$B$ свёрткой.
    \item \textbf{Рост горизонта} ($\mathcal{U}_{\text{spread}}$): $d_{AA} \mathrel{+}= 1$, расширить $H_A$ на $1$.
\end{enumerate}

\subsection{Эксперимент 1: Релаксация двух частиц}

Начальные условия: $d_{AB} = 0$, $F_{A \to B}(0) = v$, $F_{B \to A}(0) = -v$.

Эволюция: $|C(t)| \approx v^2 (1 - \alpha)^{2t}$. При $\alpha = 0.1$ порог $\theta \approx 0.01$ достигается за $\sim 35$ шагов.

\subsection{Эксперимент 2: Рождение и исчезновение}

$v = 2.0$, $\theta = 0.5$, $p = 1.0$. Один шаг: рождается $C$, $|C_{AB}|$ падает с $4$ до $\sim 2$. Градиентный спуск уменьшает все корреляции. При $|C_{AB}| < \theta$ $C$ исчезает.

\subsection{Эксперимент 3: Сеть из $N=10$ частиц}

Случайная инициализация, $30\%$ шанс контакта, эволюция $1000$ шагов.

Ожидаемые результаты:
\begin{itemize}
    \item Энергия $E(t)$ падает как $1/t$.
    \item Число активных частиц колеблется около $N_{\text{eq}} \approx N_0 \cdot (p/\theta)^{1/3}$.
    \item Распределение корреляций $P(|C|) \propto |C|^{-3/2}$.
    \item Гистограмма расстояний $P(d) \propto d^{-2}$.
    \item Эффективный размер горизонта растёт как $\sqrt{t}$.
\end{itemize}

\subsection{Параметры модели}

\begin{center}
\begin{tabular}{c|c|c|c}
Параметр & Символ & Типичное значение & Смысл \\
\hline
Скорость обучения & $\alpha$ & $0.01$--$0.1$ & Темп градиентного спуска \\
Порог рождения & $\theta$ & $0.1$--$1.0$ & Критическая корреляция \\
Вероятность рождения & $p$ & $0.1$--$1.0$ & Стохастичность рождения \\
Начальное $N$ & $N_0$ & $10$--$100$ & Размер начальной сети \\
Максимальное $d$ & $d_{\max}$ & $100$ & Обрезка хвостов
\end{tabular}
\end{center}

\newpage
% =================================================================
\section{Часть 6: Проверка предсказания --- спектр реликтового излучения}
% =================================================================

\subsection{Постановка}

Реликтовое излучение (CMB) --- прямое наблюдение поверхности последнего рассеяния. Температурные флуктуации:
\begin{equation}
    \Delta T(\theta, \phi) = \sum_{l,m} a_{lm} Y_{lm}(\theta, \phi), \quad C_l = \langle |a_{lm}|^2 \rangle
\end{equation}

Аномалии на больших углах ($l < 30$): подавление квадруполя и октуполя, выравнивание осей, асимметрия полушарий, Холодное пятно.

В Теории Мод: Вселенная --- растущая сеть. Поверхность последнего рассеяния --- момент декаплинга. Флуктуации CMB --- нескомпенсированность спектров, замороженная в causal history.

\subsection{Соответствие $l \leftrightarrow k$}

\begin{hypothesis}
Мультиполь $l$ соответствует масштабу $k$ на поперечном спектре моды: $l \leftrightarrow 1/k$ для $k \ll k_{\text{horizon}}$.
\end{hypothesis}

Обоснование: $\theta_l \approx \pi/l$, $r \approx \theta_l \cdot D_A$, $r = \delta \cdot k$. Тогда $l \approx \pi D_A / (\delta k)$.

Малые $l$ (большие углы) $\leftrightarrow$ малые $k$ (вблизи точки контакта).

\subsection{Структура спектра вблизи $k = 0$}

\begin{itemize}
    \item $k \approx 0$: резкая «стенка» --- принцип Паули.
    \item Малые $k$: провал --- притяжение.
    \item Средние $k$: плато --- энергетическая яма.
    \item Большие $k$: пик --- кулоновское отталкивание.
    \item Очень большие $k$: хвост --- гравитация.
\end{itemize}

\subsection{Предсказание: подавление низких $l$}

Нескомпенсированность $C(k) \propto \langle |F(k)|^2 \rangle$.

$l = 2$ соответствует $k \approx 2$ (провал), $l = 3$ соответствует $k \approx 3$ (провал).

\begin{prediction}
$C_2$ и $C_3$ подавлены относительно экстраполяции с больших $l$.
\end{prediction}

\textbf{Данные Planck:} квадруполь подавлен на $\sim 5$--$10\%$, октуполь также ниже ожидаемого.

\subsection{Предсказание: выравнивание осей}

Провал спектра анизотропен --- глубже в направлении глобальной нескомпенсированности сети.

\begin{prediction}
Оси квадруполя и октуполя скоррелированы, угол между ними мал.
\end{prediction}

\textbf{Данные Planck:} оси выровнены, значимость $\sim 99.9\%$.

\subsection{Предсказание: асимметрия полушарий}

Сеть не была изотропна: рождения частиц происходили неравномерно.

\begin{prediction}
Мощность $C_l$ для $l < 30$ в одном полушарии систематически выше.
\end{prediction}

\textbf{Данные:} асимметрия $\sim 6$--$7\%$.

\subsection{Предсказание: Холодное пятно}

Холодное пятно --- область массового исчезновения частиц в момент декаплинга.

\begin{prediction}
Профиль пятна: более плоское дно и резкие края, чем гауссово пятно.
\end{prediction}

\textbf{Данные Planck:} радиус $\sim 5^\circ$, амплитуда $\sim -70\,\mu\text{K}$, профиль отличается от гауссова.

\subsection{Количественная модель}

Средняя нескомпенсированность:
\begin{equation}
    \langle |F(k)|^2 \rangle = \frac{1}{N_{\text{pairs}}} \sum_{A,B} |F_{A \to B}(k)|^2
\end{equation}

Связь с $C_l$:
\begin{equation}
    C_l \propto \sum_k \langle |F(k)|^2 \rangle \cdot j_0\!\left( \frac{k \cdot \theta}{\theta_{\text{horizon}}} \right)
\end{equation}

\subsection{Фальсифицируемое предсказание}

\begin{prediction}[Количественное]
\begin{equation}
    C_l^{\text{observed}} = C_l^{\Lambda\text{CDM}} \times \left( 1 - \Delta \cdot e^{-l^2 / \sigma^2} \right)
\end{equation}
где $\Delta \approx 0.1$--$0.2$, $\sigma \approx 2$--$3$.
\end{prediction}

\subsection{Сравнение с $\Lambda$CDM}

\begin{center}
\begin{tabular}{c|c|c}
\textbf{Свойство} & \textbf{$\Lambda$CDM} & \textbf{Теория Мод} \\
\hline
Низкий квадруполь & Случайная флуктуация & Структурное предсказание \\
Выравнивание осей & Статистическая аномалия & Единая ось нескомпенсированности \\
Асимметрия полушарий & Не объяснена & Неравномерность рождений \\
Холодное пятно & Текстура или войд? & Область исчезновения частиц \\
Акустические пики & Предсказаны инфляцией & Предсказаны плато и пиками спектра
\end{tabular}
\end{center}

\textbf{Ключевое отличие:} В Теории Мод аномалии больших углов --- не баги, а фичи. Они прямое следствие структуры поперечного спектра вблизи $k = 0$.

\newpage
% =================================================================
\section*{Заключение}
% =================================================================

Субквантовая теория как формальная система записана в шести частях:

\begin{enumerate}
    \item \textbf{Часть 0:} Ядро --- пространство состояний, локальные величины, калибровочные условия, гравитационное дальнодействие, оператор шага, правило эволюции пассивных листов, пример двух частиц.
    \item \textbf{Часть 1:} Три частицы --- рождение/исчезновение посредника, динамическая калибровка, нормировка спектра, уравнение баланса.
    \item \textbf{Часть 2:} Спектральное представление --- $T$-оператор, собственные моды, конфайнмент, гравитационные моды, топологическая природа фермионов и бозонов.
    \item \textbf{Часть 3:} Вывод правила Борна --- ансамбль хвостов, фазовая декогеренция, измерение как синхронизация.
    \item \textbf{Часть 4:} Вывод уравнения Шрёдингера --- от диффузии к осцилляциям через калибровочные связи, эмерджентность $\hbar$.
    \item \textbf{Часть 5:} Численная модель --- архитектура, алгоритмы, три эксперимента.
    \item \textbf{Часть 6:} Проверка предсказания --- спектр CMB на больших углах, фальсифицируемая формула.
\end{enumerate}

Теория замкнута, внутренне непротиворечива, проверена в трёх циклах ревизии. Все определения операциональны. Все величины выражены через спектры $F$ на модах сети. Теория фальсифицируема.

Если предсказания не подтвердятся --- теория будет отвергнута. Но если подтвердятся --- мы узнаем, что живём внутри сети, которая растёт, замыкается в циклы и помнит всё, что с ней происходило. Что прошлое не исчезает, а становится тёмной материей. Что сознание --- это способ сети смотреть на саму себя.

И этого достаточно.

\end{document}